\section{The Narrow Gate and the Silence on Names}

On the road to Jerusalem, an unnamed bystander put to our Lord the question that every generation sooner or later asks: ``Lord, are they few that are saved?'' (Luke 13:23). Christ's reply is one of the most instructive moments in the Gospels. He refuses to give a census. He turns His eyes from the bystander to the whole crowd and lays the weight of the question upon the hearer's own will:

\begin{quote}
Strive to enter by the narrow gate: for many, I say to you, shall seek to enter and shall not be able. But when the master of the house shall be gone in and shall shut the door, you shall begin to stand without; and knock at the door, saying: Lord, open to us. And he answering, shall say to you: I know you not, whence you are. Then you shall begin to say: We have eaten and drunk in thy presence: and thou hast taught in our streets. And he shall say to you: I know you not, whence you are. Depart from me, all ye workers of iniquity. (Luke 13:24--27)
\end{quote}

Every word of that passage cuts through idle speculation. Familiarity with Christ---hearing His teaching in the streets, sitting near His table---saves no one who continues working iniquity. The question is not how many others will pass through the door; the question is whether \textit{you} are striving (\textit{contendite}) to enter before the Master shuts it.

Yet the Scriptures do not leave the matter in vague optimism. In the Sermon on the Mount, the Lord speaks directly of the two ways: ``Enter ye in at the narrow gate: for wide is the gate, and broad is the way that leadeth to destruction, and many there are who go in thereat. How narrow is the gate, and strait is the way that leadeth to life: and few there are that find it!'' (Matt 7:13--14). At the end of the parable of the wedding feast, when the guest without a wedding garment is bound hand and foot and cast into exterior darkness, He adds: ``For many are called, but few are chosen'' (Matt 22:14). Five of the ten virgins find the marriage door shut (Matt 25:1--12); many who cry ``Lord, Lord'' and even prophesy in His name hear Him profess, ``I never knew you'' (Matt 7:21--23).

The Fathers and medieval Doctors read these texts with unflinching sobriety. Preaching in Antioch---one of the great Christian cities of the fourth century---St.\ John Chrysostom warned his congregation in his twenty-fourth Homily on the Acts of the Apostles against taking comfort in the sheer size of the baptized multitude:

\begin{quote}
Many more in number are they who go down into hell: but greater than it is the Kingdom, however few it contain. \dots\ How many, suppose you, may there be in our city who are likely to be saved? It is disagreeable, what I am going to say, but I will say it nevertheless. Among all these myriads, there are not to be found one hundred likely to be saved: nay, even as to these, I question it. For think, what wickedness there is in the young, what supineness in the aged! \dots\ These things I say, that none may admire the Church because of its numbers, but that we may study to make the multitude proof-worthy.\endnote{St.\ John Chrysostom, \textit{Homilies on the Acts of the Apostles} 24, NPNF1 vol.\ 11.}
\end{quote}

Chrysostom speaks as a pastor watching his own flock kneel in church with ``stupidly unconcerned'' hearts, refusing to flatter them with safety in numbers.

St.\ Thomas Aquinas addresses the same question calmly and systematically in the \textit{Prima Pars} of the \textit{Summa Theologiae} (I, q.\ 23, a.\ 7), asking whether the number of the predestined is certain. In the body of the article, after mentioning various medieval speculations comparing the number of the saved to the number of the angels, Aquinas lays down a sentence drawn directly from the Roman Missal's Secret prayer \textit{pro vivis et defunctis}: ``It is, however, better to say that, `to God alone is known the number for whom is reserved eternal happiness.'\,''\endnote{St.\ Thomas Aquinas, \textit{Summa Theologiae} I, q.\ 23, a.\ 7, trans.\ Fathers of the English Dominican Province.} Yet in his Reply to the Third Objection---which had argued from Matthew 7:13--14 that more seem to be lost than saved, whereas in nature good prevails in the majority---Aquinas states his own theological judgment plainly and without evasion:

\begin{quote}
The good that is proportionate to the common state of nature is to be found in the majority; and is wanting in the minority. The good that exceeds the common state of nature is to be found in the minority, and is wanting in the majority. Thus it is clear that the majority of men have a sufficient knowledge for the guidance of life; and those who have not this knowledge are said to be half-witted or foolish; but they who attain to a profound knowledge of things intelligible are a very small minority in respect to the rest. Since [therefore] eternal happiness, consisting in the vision of God, exceeds the common state of nature, and especially in so far as this is deprived of grace through the corruption of original sin, those who are saved are in the minority. In this especially, however, appears the mercy of God, that He has chosen some for that salvation, from which very many in accordance with the common course and tendency of nature fall short.\endnote{St.\ Thomas Aquinas, \textit{Summa Theologiae} I, q.\ 23, a.\ 7, ad 3, trans.\ Fathers of the English Dominican Province (the printed English reads ``their'' where the Latin has \textit{igitur}, ``therefore'').}
\end{quote}

Honesty requires setting this common patristic and Thomistic view down without softening it. Neither Chrysostom nor Augustine nor Aquinas treated hell as an empty threat or a rare anomaly. At the same time, honesty requires setting beside it the full breadth of Catholic teaching. The Church has never defined the relative proportion of the saved and the lost, keeping intact the Missal's confession that ``to God alone is known the number for whom is reserved eternal happiness.'' Later Catholic theologians---particularly after reflection on the vast multitude of baptized infants who die before the use of reason, and on the hidden workings of grace and final purification---have differed over whether the majority of adult Christians or of mankind attain salvation. In his 2007 encyclical \textit{Spe salvi} (\S\S 45--46), Pope Benedict XVI acknowledges both poles of human destiny: while some have closed themselves utterly against truth and love so that their state becomes irrevocable, ``the great majority of people'' retain an underlying openness to truth and love beneath the compromises and stains of life and are healed through the purifying encounter with Christ in Purgatory.\endnote{Pope Benedict XVI, Encyclical \textit{Spe salvi} (30 November 2007), \S\S 45--47.} Between Aquinas's sober judgment and Benedict XVI's hopeful pastoral reflection, the Church leaves room for theological debate on the \textit{proportion}, while closing the door decisively against universalism and against any complacency that treats salvation as automatic.

Equally firm is the Church's restraint regarding individual human names. The Church canonizes saints with solemn certainty, declaring that particular men and women now behold the face of God in heaven; she has never solemnly defined of any individual human being, living or dead, that he or she is in hell. St.\ Augustine explains the root of this restraint in \textit{The City of God} XXI.24: because the Church has certainty regarding no living person's final reprobation, ``she prays for all her enemies who yet live in this world.''\endnote{St.\ Augustine, \textit{The City of God} XXI.24, trans.\ Dods.} The \textit{Catechism of the Catholic Church} sets the same boundary on human judgment: ``although we can judge that an act is in itself a grave offense, we must entrust judgment of persons to the justice and mercy of God'' (CCC 1861).

One name in Scripture stands closest to a solemn exception, and it must be handled with exact fidelity to both Gospel and liturgy: Judas Iscariot. Our Lord Himself says of His betrayer, ``woe to that man by whom the Son of man shall be betrayed. It were better for him, if that man had not been born'' (Matt 26:24), and in His High Priestly Prayer declares to the Father, ``none of them is lost, but the son of perdition: that the scripture may be fulfilled'' (John 17:12). The Roman liturgy itself speaks of Judas with astonishing candor. In the 1962 \textit{Missale Romanum}, the Collect appointed for both Holy Thursday and Good Friday sets Judas side by side with the penitent thief of Calvary (Luke 23:42--43):

\begin{quote}
\textit{Deus, a quo et Judas reatus sui poenam, et confessionis suae latro praemium sumpsit, concede nobis tuae propitiationis effectum: ut, sicut in passione sua Jesus Christus, Dominus noster, diversa utrisque intulit stipendia meritorum; ita nobis, ablato vetustatis errore, resurrectionis suae gratiam largiatur.}\endnote{1962 \textit{Missale Romanum}, Feria V in Coena Domini (Collect) and Feria VI in Parasceve (Collect after the first Lesson).}
\end{quote}

``O God, from whom both Judas received the punishment of his guilt, and the thief the reward of his confession, grant us the effect of thy propitiation: that, even as in his passion Jesus Christ our Lord rendered to each the diverse wages of his deserts, so unto us, having taken away the error of our old life, he may bestow the grace of his resurrection.''

Look closely at what that ancient Roman prayer does. It refuses to sentimentalize Judas's betrayal, holding up the real, terrifying contrast between two men who died on the same Friday---one in despair, the other in penitent confession---and both receiving the wages of their deserts from the Passion of Christ. Yet notice where the Collect turns the worshipper's gaze: not to curiosity about another man's sentence, but to his own need of propitiation, asking that our old error be stripped away and the grace of Christ's resurrection be given to us. Even when the liturgy names Judas, it speaks to keep \textit{us} awake.
